\section{Méthode d'évaluation et protocole}
\subsection{Comptages et tolérance géométrique}
Les références de Groningen ont des contours noirs sur fond blanc. Elles sont inversées par un seuillage à 127 pour obtenir des cartes binaires cohérentes avec les sorties du programme. Les mesures utilisent la résolution originale.

On note $D$ le nombre de pixels détectés, $R$ le nombre de pixels de référence et $C$ le nombre de correspondances correctes. Les mesures demandées sont :
\begin{align*}
  FP &= D-C, & FN &= R-C, & Z &= C+FP+FN,\\
  P &= \frac{C}{Z}, & TFP &= \frac{FP}{Z}, & TFN &= \frac{FN}{Z}.
\end{align*}
Ces trois taux somment à 1. Le $TFP$ du sujet diffère du taux de classification $FP/(FP+VN)$. Si les deux cartes sont vides, nous fixons $P=1$ et $TFP=TFN=0$.

En mode exact, les pixels doivent occuper la même position. En mode tolérant, une correspondance est admise dans un voisinage $3\times3$ : au plus un pixel de décalage par axe. Un \textbf{appariement maximum un à un} (Hopcroft--Karp) interdit de compter plusieurs détections pour une même référence et garantit $C\leq\min(D,R)$. Par exemple, deux détections proches d'une unique référence donnent $C=1$, $FP=1$, $FN=0$ et $P=0{,}5$.

\subsection{Choix des paramètres}
Les noms des images sont triés ; une image sur quatre est réservée au réglage, soit 10 images, et les 30 autres à l'évaluation. Nous comparons \textbf{72 configurations} sur les dix images de réglage, en maximisant la moyenne des $P$ avec tolérance $3\times3$. Le noyau est fixé à $k=3$ et $\sigma\in\{0,1,2\}$.
\begin{itemize}
  \item Sobel : $T\in\{10,20,40,80,120,160\}$ ; Laplace : $T\in\{5,10,20,40,80,120\}$.
  \item Canny : seuil bas dans $\{20,40,60,80,120,160\}$, seuil haut égal à 2 ou 3 fois le seuil bas.
\end{itemize}
Les paramètres sélectionnés sont ensuite communs à toutes les images de test, sans réglage individuel à partir de leurs références.

\begin{center}\small
\begin{tabularx}{\linewidth}{lYrr}
\toprule
\textbf{Détecteur} & \textbf{Paramètres retenus} & \textbf{$P$ réglage} & \textbf{$P$ test} \\
\midrule
Sobel & $\sigma=1$, $k=3$, $T=160$ & 0,222 & 0,178 \\
Laplace & $\sigma=2$, $k=3$, $T=20$ & 0,143 & 0,112 \\
Canny & $\sigma=2$, $k=3$, bas $=40$, haut $=120$ & 0,290 & 0,257 \\
\bottomrule
\end{tabularx}
\captionof{table}{Paramètres figés et performances moyennes avec tolérance $3\times3$.}
\end{center}
